\begin{afterword}
\summaryhead

The author strongly suspects that a better art of rationality is possible, and says that what a group can do about it depends overwhelmingly on ways to verify which ideas work. The post sorts such methods into three levels.

The reputational level is the martial-arts master who fights realistic duels against other schools. Schools that do not compete at all end up with masters who believe in chi powers, like the schools of psychoanalysis. The experimental level compares randomly assigned teaching methods with standard measurements. Happiness research shows that a simple self-report can work when averaged over a hundred people. The organizational level needs tests of individuals that resist gaming, since a gameable test, like the SAT or a hedge fund's returns, makes a whole field adapt to the test and lose its purpose. Noisy measures can serve experiments but not judgments of individuals.

The post ends by asking readers to brainstorm measures at all three levels, even stupid ones.

\responsehead

The post is a request for ideas, and it should be judged as one. It sets out a frame and asks readers to fill it. The frame is clear, and its two technical points are correct. Noise in a measure can be averaged away in an experiment but not in a judgment of one person. A test that controls rewards invites people to game it.

Both points have earlier statements that a reader may find useful. Campbell's law, from 1976, says that when test scores become the goal of teaching they lose their value as a measure and distort the teaching. The hedge-fund example has a well-known illustration in Andrew Lo's hypothetical fund, which sold put options every month and showed returns far above the market's in a record that did not reveal its exposure to a crash. The post's model of a working measure, the self-reported happiness scale, is supported by validation studies such as that of Sandvik, Diener and Seidlitz.

The claim the post argues least is its thesis. What a group can do about rationality depends ``overwhelmingly,'' it says, on methods of verification. The martial arts and psychoanalysis show that verification matters. They do not compare it with other obstacles, such as the difficulty of working in groups, which ``A Sense That More Is Possible'' listed among the reasons there is no art of rationality, and which the next post in the book, ``Why Our Kind Can't Cooperate,'' takes up.

For information, some tools of the kind the post asks for already existed. The index of Adult Decision-Making Competence, published in 2007, measured individuals on seven decision tasks, and higher scorers reported fewer bad outcomes in their lives. Later work shows how hard the question has proved. Stanovich, West and Toplak built a comprehensive test of rational thinking, published in 2016. The Center for Applied Rationality's 2015 study of its own workshops found gains in self-reports but no change on its measures of cognitive biases, in a study without a control group that its author judged underpowered for those measures.

In short: a clear frame for testing rationality training, with correct points on noise and gaming that match earlier work such as Campbell's law; its thesis that verification matters ``overwhelmingly'' is asserted, and, as a call for ideas, the post proposes no test of its own.
\end{afterword}
